\chapter{두 지도를 섞으면?}

지금 우리에게는 지도가 두 장 있어요.
허락 지도에서 놀면 엔도스랭크, 송금 지도에서 놀면 송금 랭크예요.
그렇다면 자연스럽게 이런 생각이 들어요.
\emph{두 지도를 함께 쓰면 둘의 좋은 점을 다 가질 수 있지 않을까?}

저도 그렇게 생각했어요.
그래서 두 가지 섞는 방법을 만들어 보았어요.

\section{동전 던지기: 동전 랭크}

교실 바닥에 두 지도를 겹쳐 놓았다고 상상해 보세요.
파란 허락 화살표와 주황 송금 화살표가 함께 그려져 있어요.

쪽지 놀이 규칙을 하나만 바꿔요.
“화살표를 따라가” 카드를 뽑으면, 동전을 한 번 더 던져요.

\begin{itemize}
\item 앞면이 나오면 \textcolor{allowcol}{\textbf{파란 허락 화살표}} 가운데 하나를 따라가요.
\item 뒷면이 나오면 \textcolor{sendcol}{\textbf{주황 송금 화살표}} 가운데 하나를 따라가요.
\end{itemize}

동전을 던지니까 이 방법을 \textbf{동전 랭크}라고 부를게요.
논문에서는 “두 지도를 짝지어 묶은 페이지랭크”라는 뜻의 영어 머리글자를 따서 C-PR이라고 불러요.

동전 랭크의 두 지도는 엔도스랭크와 송금 랭크의 지도와 화살표는 같지만, 무게를 적는 법이 달라요.
허락한 양과 보낸 양을 “한 번”으로 줄이지 않고 그대로 무게로 적어요.
송금 화살표에는 신선도를 곱하고, 허락 화살표에는 곱하지 않아요.
선생님은 쪽지를 모두에게 똑같이 줘요.
이 규칙들은 동전 랭크를 처음 정한 9월 14일에 함께 정해 두었어요.

동전이 앞면일 확률을 그리스 문자 $\lambda$(람다)로 써요.
보통 동전이면 $\lambda = 0.5$예요.
$\lambda = 1$이면 늘 앞면이 나와서 허락 화살표만 따라가요. 이것을 \textbf{앞면 랭크}라고 부를게요.
$\lambda = 0$이면 늘 뒷면이 나와서 송금 화살표만 따라가요. 이것은 \textbf{뒷면 랭크}예요.
앞면 랭크는 엔도스랭크와 같은 허락 지도를 걷지만 화살표의 무게가 달라요.
뒷면 랭크도 송금 랭크와 같은 송금 지도를 걷지만 무게와 선생님 규칙이 달라요.
그러니까 동전 랭크는 앞면 랭크와 뒷면 랭크 사이 어딘가에 있는 방법이에요.
저는 $\lambda$를 0.5로 정하고, 0.25와 0.75도 함께 시험했어요.

\begin{figure}[h]
\centering
\begin{tikzpicture}[scale=0.9, every node/.style={transform shape}]
  \node[kid] (h) at (0,0) {하준};
  \node[kid] (s) at (3.6,1.3) {서연};
  \node[robot] (r) at (3.6,-1.3) {로봇};
  \draw[allow] (h) -- node[above, sloped, font=\scriptsize, color=black] {앞면이면} (s);
  \draw[send] (h) -- node[below, sloped, font=\scriptsize, color=black] {뒷면이면} (r);
  \node[draw=sunline, fill=sunfill, circle, minimum size=9mm, font=\small] at (-1.9,0) {동전};
\end{tikzpicture}
\caption{동전 랭크에서는 쪽지를 넘길 때마다 동전을 던져 어느 지도의 화살표를 따라갈지 정해요.}
\end{figure}

\section{출발점 바꾸기: 씨앗 랭크}

두 번째 방법은 조금 달라요.
쪽지는 송금 화살표만 따라가요.
그 대신 “선생님께” 카드가 나왔을 때 선생님이 쪽지를 아무에게나 주지 않고,
\emph{앞면 랭크 점수가 높은 친구}에게 더 자주 줘요.
믿음을 많이 받은 친구에게서 출발해서 송금 길을 걷는 셈이에요.
허락 지도가 뿌린 씨앗에서 자라난다는 뜻으로 이 방법을 \textbf{씨앗 랭크}라고 부를게요.
논문에서는 S-PR이라고 불러요.

\section{결과를 보기 전에 정한 약속}

새 방법을 만들면 자꾸 좋은 점만 찾고 싶어져요.
그래서 저는 동전 랭크 점수를 하나도 계산하기 전인 2026년 9월 14일에,
동전 랭크가 “성공”이라고 말하려면 무엇을 넘어야 하는지 미리 적어 두었어요.
지도교수님께 보낸 계획서와 프로그램 설정 파일에 함께 적었어요.

\begin{center}
\begin{tcolorbox}[colback=white, colframe=inkblue, boxrule=0.8pt, arc=2mm, width=0.92\linewidth]
\small
\textbf{9월 14일의 약속} \quad ($\lambda = 0.5$인 동전 랭크에 대해)
\begin{enumerate}
\item 새 허락을 맞히는 시험에서 앞면 랭크보다 못하지 않을 것.
\item 새 송금자를 맞히는 시험에서 뒷면 랭크보다 나을 것.
\end{enumerate}
두 가지를 모두 넘어야 성공이에요.
\end{tcolorbox}
\end{center}

동전 랭크가 자기를 이루는 두 지도 가운데 한쪽만 걷는 것보다 나은지를 묻는 약속이에요.
“새 허락을 맞히는 시험”과 “새 송금자를 맞히는 시험”이 무엇인지는 13장에서 자세히 설명할게요.
지금은 이것만 기억해 주세요.
\emph{결과를 보기 전에 약속을 먼저 적었다.}
이 약속이 이 책의 끝부분에서 아주 중요해져요.

\begin{deeper}[진짜 정의]
동전 랭크는 점마다 두 층의 이동 확률을 섞어요. 점 $i$에서 $j$로 갈 확률은
\[
P_{ij} = \lambda\, P^{\text{허락}}_{ij} + (1-\lambda)\, P^{\text{송금}}_{ij}
\]
이고, 한 층에 화살표가 없는 점은 다른 층의 화살표를 써요.
허락 층의 무게는 최신 허락 값의 합, 송금 층의 무게는 보낸 양에 신선도를 곱한 값의 합이에요.
논문에서는 앞면 랭크와 뒷면 랭크를 C-PR ($\lambda=1$), C-PR ($\lambda=0$)으로 써요.
9월 14일 규칙과 14장의 등록 문서는 이 둘을 EndorseRank와 AWP라고 불러요.
씨앗 랭크는 송금 층의 페이지랭크인데, 되돌림(선생님 몫)을 앞면 랭크 점수에 비례해 나누어요.
앞면 랭크 점수가 0인 점은 되돌림을 받지 않아요.
9월 14일에는 이 규칙과 함께, 같은 기간 시험에서의 규칙(이차 규칙)도 정해 두었어요.
\end{deeper}

\begin{learned}
\begin{itemize}
\item 동전 랭크는 쪽지를 넘길 때마다 동전을 던져 허락 지도와 송금 지도 가운데 하나를 고르는 방법이에요. 두 지도에는 양을 그대로 무게로 적어요.
\item 늘 앞면이면 허락 지도만 걷는 앞면 랭크, 늘 뒷면이면 송금 지도만 걷는 뒷면 랭크예요.
\item 씨앗 랭크는 앞면 랭크가 좋아하는 친구에게서 출발해 송금 지도를 걷는 방법이에요.
\item 동전 랭크가 성공했다고 말하는 기준은 결과를 보기 전인 9월 14일에 미리 적어 두었어요.
\end{itemize}
\end{learned}
